\section{Limitations}
\label{sec:discussion}
\capg{} preserves its protection under purification, and its limitations follow from how it does so. Preservation does not raise the level of protection: without purification, the edit failure rate of \capg{} is not higher than that of the additive methods and is lower than that of PhotoGuard and Mist (\cref{tab:isr}), and this rate has so far been measured on InstructPix2Pix only. The property that carries the perturbation through purification, a global mapping of colors, also determines which operations could undo it. Color-correction operations such as automatic white balance or color normalization in Lab space estimate and invert a global color mapping; they are not part of the main evaluation and may remove the protection. The same mapping leaves a visible color shift in the protected image, which affects uses that rely on faithful color, such as forensic or medical imaging. Because the defense requires only the image, a party that does not own it could also apply it to hinder legitimate edits. We therefore intend \capg{} to be applied by image owners before release, and we report the full evaluation protocol so that later attacks and defenses can be compared on the same terms.
